\section{The active best-seed branch as a one-corner odometer}

In this section we rewrite the current $d=5$ theorem route in its cleanest
form. The theorem object remains minimal:
\[
B=(s,u,v,\lambda,\mathrm{family}),\qquad \tau,\qquad \epsilon_4.
\]
Auxiliary coordinates such as $\rho$, $\alpha$, $\delta$, $\beta$, and
\[
\kappa:=q+s+v+\lambda \pmod m
\]
are used only for proof or constructive purposes and are not part of the main
theorem statement.

Fix an odd modulus $m\ge 5$ and the best seed
\[
L=[2,2,1],\qquad R=[1,4,4].
\]
We work on the active nonterminal branch of the unresolved channel
$R1\to H_{L1}$.

\begin{proposition}[explicit trigger family]
The target hole family is
\[
H_{L1}=\{(q,w,u,\lambda)=(m-1,m-1,u,2):u\neq 2\}.
\]
\end{proposition}

\begin{proposition}[universal first-exit targets]
The active branch first exits to $H_{L1}$ at exactly one of the two states
\[
T_{\mathrm{reg}}=(m-1,m-2,1),
\qquad
T_{\mathrm{exc}}=(m-2,m-1,1),
\]
with regular source families exiting at $T_{\mathrm{reg}}$ and the exceptional
source family exiting at $T_{\mathrm{exc}}$.
\end{proposition}

\begin{proposition}[pre-exit $B$-region invariance]
Every active nonterminal current state is $B$-labelled.
\end{proposition}

\begin{theorem}[phase-corner theorem]
On the active nonterminal branch, let
\[
\kappa=q+s+v+\lambda \pmod m,
\qquad
c=\mathbf 1_{\{q=m-1\}}.
\]
Then:
\begin{enumerate}
\item $\kappa'=\kappa+1\pmod m$.
\item The current event $\epsilon_4$ is determined by $\kappa$, $c$, and the
single corner condition $s=2$:
\[
\kappa=0\Rightarrow \epsilon_4=\mathrm{wrap},
\qquad
\kappa=1\Rightarrow \epsilon_4=\mathrm{carry\_jump},
\]
\[
\kappa=2\Rightarrow
\begin{cases}
\epsilon_4=\mathrm{other}_{1000},& c=1,\\
\epsilon_4=\mathrm{flat},& c=0,
\end{cases}
\]
\[
\kappa=3\Rightarrow
\begin{cases}
\epsilon_4=\mathrm{other}_{0010},& (c=0 \text{ and } s=2) \text{ or } (c=1 \text{ and } s\neq 2),\\
\epsilon_4=\mathrm{flat},& \text{otherwise},
\end{cases}
\]
and for $4\le \kappa\le m-1$ one has $\epsilon_4=\mathrm{flat}$.
\item Among flat states, the only short reset occurs at the corner
\[
(\kappa,s)=(2,2).
\]
\end{enumerate}
Equivalently, the active branch is an odometer with one corner.
\end{theorem}

\begin{proof}
By the pre-exit $B$-region invariance proposition, every active nonterminal
current state is $B$-labelled, so the outgoing direction is exactly the
unmodified mixed witness direction.

For color $0$, the mixed witness uses anchor $1$ at layer bucket $0$, anchor
$4$ at layer bucket $1$, the old-bit orientation $0/2$ at layer bucket $2$
with old bit $c=\mathbf 1_{\{q=m-1\}}$, the predecessor switch at layer bucket
$3$, and anchor $0$ at every layer bucket $4,\dots,m-1$.
Because
\[
\kappa=q+s+v+\lambda=q+w+u+v+\lambda=\sum_{i=0}^4 x_i \pmod m,
\]
this layer bucket is exactly the current value of $\kappa$. Every anchor
increments the total coordinate sum by one, hence
\[
\kappa'=\kappa+1\pmod m.
\]

Now translate anchors into grouped events. Anchor $1$ gives wrap; anchor $4$
gives carry\_jump; at $\kappa=2$ the old-bit rule chooses anchor $2$ exactly
when $c=1$, giving $\mathrm{other}_{1000}$, and chooses anchor $0$ otherwise,
giving flat. At $\kappa=3$, the predecessor flag is equivalent to $s=2$ on the
active branch, so anchor $3$ is chosen exactly in the two cases
$(c=0,s=2)$ and $(c=1,s\neq 2)$, giving $\mathrm{other}_{0010}$, while anchor
$0$ is chosen otherwise, giving flat. For $\kappa\ge 4$, anchor $0$ is used, so
the event is always flat.

Thus the stated scheduler holds. The only short flat reset is the one seen when
this scheduler reaches the predecessor-triggered branch at $(\kappa,s)=(2,2)$,
so the branch is an odometer with one corner.\qedhere
\end{proof}

\begin{corollary}[countdown law]
Define $\tau$ by $\tau=0$ on nonflat states and, on flat states,
\[
\tau=
\begin{cases}
1,& (\kappa,s)=(2,2),\\
m-\kappa,& \text{otherwise}.
\end{cases}
\]
Then
\[
\tau' = \tau-1 \qquad (\tau>0).
\]
\end{corollary}

\begin{corollary}[CJ reset law]
On a carry-jump state,
\[
R_{\mathrm{cj}}=
\begin{cases}
0,& s+v+\lambda\equiv 2 \pmod m,\\
1,& s=1 \text{ and } s+v+\lambda\not\equiv 2 \pmod m,\\
m-2,& \text{otherwise.}
\end{cases}
\]
\end{corollary}

\begin{proof}
A carry-jump state has $\kappa=1$, hence
\[
q\equiv 1-s-v-\lambda \pmod m.
\]
The zero-reset fiber is exactly $q=m-1$, equivalently
$s+v+\lambda\equiv 2\pmod m$. If $q\neq m-1$, then the successor is flat at
phase $\kappa=2$. The only short flat reset occurs at the corner
$(\kappa,s)=(2,2)$, which here corresponds to $s=1$ before the jump. Hence the
reset is $1$ when $s=1$, and $m-2$ otherwise.
\end{proof}

\begin{corollary}[OTH reset laws]
One has
\[
(1,0,0,0)\mapsto
\begin{cases}
m-3,& s=1,\\
0,& s\neq 1,
\end{cases}
\qquad
(0,0,1,0)\mapsto m-4.
\]
\end{corollary}

\begin{proof}
The event $\mathrm{other}_{1000}$ occurs at $\kappa=2$ with $c=1$. Its
successor lies at $\kappa=3$. If $s=1$, the successor is flat and the next
nonflat event occurs after $m-3$ steps; otherwise the successor is already on a
reset boundary, giving value $0$.

The event $\mathrm{other}_{0010}$ occurs at $\kappa=3$. Its successor lies at
$\kappa=4$, where the scheduler is always flat, so the next nonflat event
occurs after exactly $m-4$ steps.
\end{proof}

\begin{corollary}[finite-cover compatibility]
The phase-corner machine is compatible with the finite-cover chain
\[
B \leftarrow B+c \leftarrow B+c+d,
\]
where
\[
c=\mathbf 1_{\{q=m-1\}},
\qquad
d=\mathbf 1_{\{\text{next carry }u\ge m-3\}}.
\]
\end{corollary}

\begin{remark}[constructive cyclic controller]
The constructive route may introduce
\[
\rho=u_{\mathrm{source}}+1,
\qquad
\alpha=\rho-u,
\qquad
\beta=\alpha-(s+v+\lambda)-J,
\]
where $J=\mathbf 1_{\{\epsilon_4=\mathrm{carry\_jump}\}}$. On the checked range,
$\beta$ has unit drift and determines $q$, $c$, $\tau$, and $\mathrm{next}\,\tau$
from $(B,\beta)$. This is a constructive refinement, not a change in the main
theorem object.
\end{remark}
